\clearpage
\VerdierPhysicalPageAnchor{0239}%
\begingroup\small\setlength{\parskip}{0pt}\setlength{\emergencystretch}{1.5em}
\setlength{\abovedisplayskip}{4pt plus 2pt minus 2pt}%
\setlength{\belowdisplayskip}{4pt plus 2pt minus 2pt}%

\VerdierSemanticAnchor{ch4d-4-4-2-b-cont}
\noindent\textit{One then has}
\[
\mathrm I_2^{p,q}=H_{\mathcal B}^pF^{q,\bullet}(Y^\bullet).
\]

\VerdierSemanticAnchor{ch4d-4-4-2-c}
\noindent\textit{c) The terms $\mathrm I_\infty^{p,q}$ are the graded
objects associated with the suitably filtered objects $F^{p+q}(Y)$.}

\VerdierSemanticAnchor{ch4d-4-4-2-proof}
We prove a). Let $\bigl(Y(p,q),\delta(p,q,r)\bigr)$ be the first spectral
object with values in $K(\mathcal A)$ associated with the complex $Y$. It
follows immediately from the definition of this object that (4.1.2.3)
gives
\[
Y(-p-1,-p)=Y^p[-p].
\]
Consequently (Chapter~II, 4.3.3.1),
\[
\mathrm I_1^{p,q}=F^{p+q}\bigl(Y^p[-p]\bigr)=F^q(Y^p).
\]
The description of the differential
$d_1\colon\mathrm I_1^{p,q}\longrightarrow\mathrm I_1^{p+1,q}$ now
follows immediately from (4.1.2.6). Assertion b) follows immediately from
a). The last assertion follows from (Chapter~II, 4.3.2).

\VerdierSemanticAnchor{ch4d-4-4-3}
\medskip\noindent\textbf{Remark 4.4.3.}\quad
Denote by $\bigl(\mathrm I_r^{p,q}(Y),d_r,r\geq1\bigr)$ the first spectral
sequence of the functor $F$ relative to the complex $Y$. The functor
\[
Y\longmapsto\bigl(\mathrm I_r^{p,q}(Y),d_r,r\geq1\bigr)
\]
is a functor from the category $\operatorname{comp}(\mathcal A)$ to the
category of spectral sequences with values in $\mathcal B$. This functor
does not in general factor through the category $K(\mathcal A)$. On the
other hand, the functor

\VerdierSemanticAnchor{ch4d-4-4-3-r-ge-2}
\[
Y\longmapsto\bigl(\mathrm I_r^{p,q}(Y),d_r,r\geq2\bigr)
\]
factors, by Proposition (4.4.2), b), through the category of complexes up
to homotopy.

\VerdierSemanticAnchor{ch4d-4-4-4}
\medskip\noindent\textbf{4.4.4.}\quad
Now let $\mathcal A$ be an abelian category, let
$F\colon D(\mathcal A)\longrightarrow\mathcal B$ be a cohomological
functor with values in an abelian category $\mathcal B$, and let $F^p$,
$p\in\mathbb Z$, be the functors obtained from $F$ by composition with the
translations of $D(\mathcal A)$ ($F^p(X)=F(X[p])$). In addition, let $Y$
be an object of $\operatorname{comp}(\mathcal A)$, and let
$\bigl(Y(p,q),\delta(p,q,r)\bigr)$ be the first spectral object, with values
in $D(\mathcal A)$, associated with the complex $Y$. The functor $F$ takes
this spectral object to a spectral object with values in $\mathcal B$.
This yields a spectral sequence with values in $\mathcal B$, again called
the \emph{first spectral sequence of the functor $F$ relative to the
complex $Y$}. This spectral sequence is described essentially by
Proposition (4.4.2) (in Proposition (4.4.2), it suffices to replace the
category $K(\mathcal A)$ by the category $D(\mathcal A)$). The first
spectral sequence of the functor $F$, relative to the complex $Y$, depends
functorially on $F$ and on the argument $Y$ when the latter varies in
$\operatorname{comp}(\mathcal A)$. For this spectral sequence we shall
always use the indexing introduced in (Chapter~II, 4.3.3.1).
\endgroup
